\chapter{总结与展望}
本章对面向虚拟现实动画网格的时空感知注视点渲染技术研究的研究工作进行总结，并对未来工作进行展望。

\section{论文总结}

本文聚焦于人类视觉系统的时空感知特性，研究了面向虚拟现实动画网格的时空感知注视点渲染技术，通过有效利用扫视过程中的时序感知冗余与动画网格的时空相关性，在保证视觉感知质量的前提下显著提升渲染性能。具体研究内容如下：

（1）针对传统注视点渲染存在的多维感知失配与时序感知冗余问题，提出融合扫视时序特征的多维感知注视点渲染方法。首先，通过显式引入扫视时序调制因子，综合集成色彩、刺激面积、时空频率、亮度、视网膜偏心度等多维参数，构建了融合扫视时序特征的多维视觉感知模型castleCSF-TS，精确量化了扫视后视觉敏感度动态恢复过程。其次，提出感知自适应的动态分辨率渲染算法，通过高精度眼动事件检测划分时序优化窗口，并依据castleCSF-TS模型实时预测的感知临界频率，在时序优化窗口内动态调整中央凹及周边区域的渲染分辨率，使渲染负载与瞬时视觉感知能力的精准对齐，从而显著降低时序计算冗余，提升渲染效率。最后，实验结果表明，本文方法在高动态交互场景中，在保持视觉感知质量的同时有效提升了渲染性能与系统能效。

（2）针对传统LOD方法在处理动画网格时存在的运动失真与时空资源错配问题，提出基于时空视觉感知的动画网格混合LOD方法。首先，在离线预处理阶段提出时空运动特征驱动的动画网格层次简化算法，通过构建融合动画时序特性与空间结构特征的重要性度量模型，指导生成在动画播放全周期内保持视觉连贯性的多层次细节表示。在实时渲染阶段提出感知引导的自适应细分LOD算法，利用 GPU 硬件镶嵌单元，依据castleCSF-TS 模型计算得到的分辨率缩放因子，在基础 LOD 层级上进行像素级精度的表面细分，实现了动画网格在时空维度上的平滑细节过渡和高效资源分配。最后，实验结果表明，本文方法能够在显著降低几何渲染开销的同时，有效保障动态场景中视觉体验的连贯性与真实感。

本文所提出的融合扫视时序特征的多维感知注视点渲染方法和基于时空视觉感知的动画网格混合LOD方法为改进传统注视点渲染技术提供了有效的理论支撑，将上述方法应用于实际生产中，设计并实现了能够以第一人称视角漫游检索并进行数据分析的数字工厂可视化系统。通过功能验证和性能分析，所提方案有效提高了系统的整体性能、可靠性和交互性，降低了对硬件设备的高要求，提升了用户的使用体验。同时，该技术也
适用于其它虚拟现实应用。

\section{未来工作与展望}
尽管本文在动画网格的时空感知渲染方面取得了一定的研究成果，但结合当前感知研究的前沿趋势及复杂交互环境的需求，仍存在以下方向值得进一步深入探讨：

（1）注视点预测机制与系统延迟补偿的优化。本文虽然引入了预测机制，但受限于硬件传输延迟，在极快速扫视下仍存在感知的微小波动。未来工作可借鉴更先进的注视点轨迹预测模型，探索在渲染管线前端进行前瞻性资源预分配，利用预测的注视着陆点提前执行低分辨率渲染，以有效补偿系统延迟带来的影响。

（2）复杂光影与全色域环境下的时空感知泛化研究。目前castleCSF-TS 模型主要验证于几何与着色频率维度。未来工作可进一步探索在高对比度场景和宽色域显示设备下，扫视运动对亮度掩蔽和色彩空间对比敏感度的耦合影响，将 castleCSF-TS 模型扩展至全色域感知空间，以适配下一代高规格 VR 硬件。

（3）动画网格拓扑一致性与感知伪影的深度融合 在极高倍率的网格简化下，边缘区域的几何突变易导致视觉闪烁或走样。未来工作可将时域抗锯齿（TAA）技术与感知模型进行深度耦合，根据 castleCSF-TS 模型动态调整 TAA 的时域混合权重。利用历史帧信息补偿眼跳期间或低分辨率区域丢失的几何细节，通过时域平滑机制消除 LOD 切换带来的伪影，在追求极致压缩率的同时确保视觉连贯性。